\appendix
\section{The low-rank cohomology and the building}
\label{app:background}

We give the two elementary ingredients in the proof of
\Cref{prop:group-classes}.  The calculations are over characteristic
zero; no reduction modulo $p$ of the matrices below is used.
The degree-three trace class is the classical cocycle associated to an
invariant symmetric bilinear form
\cite[\S21]{ChevalleyEilenberg1948}, \cite[\S11]{Koszul1950}.
We give the low-rank calculation and block restriction in the
normalization used in this paper.

\subsection{A finite Chevalley--Eilenberg calculation}
\label{app:ce}

For $\mathfrak{sl}_n$, order the matrix units $E_{ij}$, $i\ne j$,
lexicographically, followed by $H_i=E_{ii}-E_{nn}$ for $1\le i<n$.
Order exterior bases by increasing tuples, lexicographically, starting
with index zero.  The differential is determined by
\[
 (d\omega)(X_0,\ldots,X_k)=
 \sum_{i<j}(-1)^{i+j}
 \omega([X_i,X_j],X_0,\ldots,\widehat X_i,\ldots,
                         \widehat X_j,\ldots,X_k).
\]
For $n=2$, the matrix of $d_1$ in these bases is
\[
 \begin{pmatrix}0&0&-1\\2&0&0\\0&-2&0\end{pmatrix},
\]
with determinant $4$, and $d_2=0$.  For $n=3$, the dimensions of
$C^0,\ldots,C^4$ are $1,8,28,56,70$.  The following nonsingular
minors of $d_1,d_2,d_3$ provide a short exact-arithmetic certificate.
Their row and column indices refer to the bases just specified:
\begin{align*}
 R_1={}&(1,2,3,5,7,9,15,16),&
 C_1={}&(0,1,2,3,4,5,6,7),\\
 \det(d_1[R_1,C_1])={}&-1.&&
\end{align*}
For the other two minors, use
\begin{align*}
 R_2={}&(0,2,4,6,7,8,9,11,12,13,19,22,24,28,30,37,39,41,43,52),\\
 C_2={}&(0,1,2,3,4,5,6,7,8,11,13,14,15,16,17,20,22,23,25,27),\\
 \det(d_2[R_2,C_2])={}&2916,
\end{align*}
\begin{align*}
 R_3={}&(0,1,2,3,4,5,6,7,9,10,11,12,15,16,17,18,19,\\
       &\hspace{1em}20,21,22,23,25,26,27,31,35,36,37,39,40,41,45,55,56,61),\\
 C_3={}&(0,1,2,3,5,6,7,8,9,10,11,12,13,14,15,16,17,18,\\
       &\hspace{1em}20,21,22,24,25,28,31,32,36,37,38,40,42,43,44,45,52),\\
 \det(d_3[R_3,C_3])={}&1024.
\end{align*}
These integer identities can be verified directly from the displayed
differential and the commutator formula
\[
 [E_{ij},E_{kl}]=\delta_{jk}E_{il}-\delta_{li}E_{kj}.
\]
The script \texttt{support/r-ce/ce\_certificate.py} generates these
integer matrices and evaluates the specified minors by fraction-free
elimination, using Python's standard library. To turn these rank bounds
into cohomology dimensions, we now prove that the trace cocycle is not
a boundary.

The alternating cochain $\omega(X,Y,Z)=\tr(X[Y,Z])$ is closed by
trace invariance and the Jacobi identity.  Its restriction to the
$\mathfrak{sl}_2$ block satisfies
$\omega(E_{12},E_{21},E_{11}-E_{22})=2$.  On this block every
three-coboundary is zero: substituting these three vectors in $d\eta$
gives
$-\eta(h,h)+\eta(-2e,f)-\eta(2f,e)=0$.
Thus the trace cochain is not a boundary, in either rank.
The first minor gives $\rank d_1=8$ for $\mathfrak{sl}_3$.
Since $d^2=0$, the second rank is at most $28-8=20$; its minor gives
equality.  The kernel of $d_3$ contains this twenty-dimensional image
and the independent trace class, so its rank is at most $56-21=35$;
the last minor gives equality.  Therefore, for $n=2,3$,
\[
 H^1(\mathfrak{sl}_n)=H^2(\mathfrak{sl}_n)=0,
 \qquad H^3(\mathfrak{sl}_n)=\Q_p\cdot[\omega].
\]
The direct sum $\mathfrak{gl}_n=\mathfrak{sl}_n\oplus\Q_p$ gives
exactly the low-degree groups used in \Cref{prop:group-classes}.
Block restriction of $\omega$ is the corresponding trace cochain.
For a central simple inner form, extend scalars to a finite splitting
field and use reduced trace.  The finite-dimensional cochain complex
commutes with this faithfully flat extension, so the calculation descends.

We also need inner invariance in every degree for the cellular argument.
Cartan's identity $\mathcal L_X=d\iota_X+\iota_Xd$ shows that the Lie
algebra of inner automorphisms acts trivially on cohomology.  After
splitting, the action is algebraic and the group $\GL_n$ is connected
in characteristic zero; hence its action on cohomology is trivial.
Faithfully flat descent gives the same assertion for the inner forms.

\subsection{An elementary model of the building}
\label{app:building}

Let $V=\Q_p^n$, $\cO=\Z_p$, and $|p|=p^{-1}$.  The vertices of
$\mathcal B_n$ are homothety classes of $\cO$-lattices.  A simplex is a
chain of distinct lattices between $L$ and $pL$, modulo homothety.
A maximal chain corresponds to a complete flag in $L/pL$, so its
realization has dimension $n-1$.

Here is a norm description, including contractibility.  Every
nonarchimedean norm $\alpha$ on $V$ has an orthogonal basis.
Indeed, for a fixed maximum norm $N$, the triangle inequality gives
$\alpha\le CN$.  Continuity and compactness of $N(v)=1$ give a positive
lower bound $cN\le\alpha$.  Thus every positive-radius closed ball is
a lattice.  Put $L=\{\alpha\le1\}$, so $pL=\{\alpha\le p^{-1}\}$.
On the compact set $L\setminus pL$ the norm is locally constant and
has finitely many values in $(p^{-1},1]$.  Its balls modulo $pL$ form
a flag.  Lift an adapted residue basis, choosing each lift in the
corresponding ball.  For a linear combination with largest coefficient
norm one, use the largest basis weight among the unit coefficients.
Its nonzero residue in that flag quotient shows that the sum has
exactly this norm; nonunit coefficients contribute at most $p^{-1}$.  Scaling
proves orthogonality.

Norms modulo positive real scaling identify with $|\mathcal B_n|$.
Explicitly, for an adapted basis $e_i$ put
$L_j=p\cO e_1+\cdots+p\cO e_j+\cO e_{j+1}+\cdots+\cO e_n$.
Barycentric coordinates $(\theta_0,\ldots,\theta_{n-1})$ on this chamber
give the split norm with weights $p^{a_i}$, where
\[
 a_i=\sum_{j=i}^{n-1}\theta_j\quad(i<n),\qquad a_n=0.
\]
Conversely $\theta_0=1-a_1$ and $\theta_j=a_j-a_{j+1}$.
The periodic chain of norm balls and the ratios between its jumps recover
these data.  Vanishing barycentric coordinates remove precisely the
corresponding jumps, and the endpoint identification is multiplication
of the lattice by $p$.  This proves the claimed identification on faces
as well as chambers.

For two norms let $M(\alpha,\beta)$ and $m(\alpha,\beta)$ be the
maximum and minimum of $\alpha(v)/\beta(v)$ on $V\setminus0$.
They exist by compactness of $\mathbf P(V)$.  The metric
$d([\alpha],[\beta])=\log M-\log m$ induces the simplicial topology.
To verify the only local issue, normalize $m(\alpha,\beta)=1$.
A bounded metric neighborhood satisfies
$\beta\le\alpha\le e^R\beta$.  For $p^{-1}\le r\le1$ all its
lattices $\{\alpha\le r\}$ lie between the two fixed lattices
$\{\beta\le p^{-1}e^{-R}\}$ and $\{\beta\le1\}$.
There are finitely many such lattices; this radius interval represents
every ball homothety class.  Thus the neighborhood meets finitely many
simplices.  On each simplex the weight formula is continuous, and on
a finite union of closed simplices compactness gives the inverse
continuity.  This proves the topology assertion.

With the same normalization define
\[
 H_t(\alpha)(v)=\max\{t\alpha(v),\beta(v)\},\qquad 0\le t\le1.
\]
This is a norm, remains normalized, equals $\beta$ at $t=0$ and
$\alpha$ at $t=1$, and fixes $\beta$.  In logarithmic ratios it is the
map $f\mapsto\max\{\log t+f,0\}$ for $t>0$.
Metric convergence of normalized norms is equivalent to uniform
convergence of these ratios.  Near any fixed norm they are uniformly
bounded, so for sufficiently small $t$ the displayed maximum is
identically zero.  This also proves continuity at $t=0$ and hence
contractibility.

Finally, the type of $[g\cO^n]$ is $v_p(\det g)$ modulo $n$.
Each chamber has every type once and $J_n$ preserves types.
Start a target chamber at its unique type-zero vertex.  Scaling its
lattice by an integral power of $p$ makes its determinant valuation zero;
an adapted basis for its residue flag then carries the standard chamber
to it by an element of $J_n$.  Thus $J_n$ is chamber-transitive.
A vertex stabilizer in $J_n$ actually preserves its lattice, since
$gL=aL$ implies $0=v_p(\det g)=n v_p(a)$.  It is consequently a
conjugate of $\GL_n(\Z_p)$, compact open.  Face stabilizers are finite
intersections of these groups and fix their vertices pointwise because
types are distinct.  The closed chamber is therefore a fundamental
domain, including faces, with no inversions.  These are exactly the
properties used by the finite cellular resolution in \cref{sec:cohomology}.
